\chapter{Additional Derivations}
\label{app:derivations}

\section{Merging two processed key sets}

Let $J_1$ and $J_2$ be disjoint processed key sets for one query row. Define
\begin{align}
  m_a &= \max_{j\in J_a} s_{ij}, \\
  \ell_a &= \sum_{j\in J_a} e^{s_{ij}-m_a}, \\
  A_a &= \sum_{j\in J_a} e^{s_{ij}-m_a} v_j,
\end{align}
for $a\in\{1,2\}$. Let $m = \max(m_1,m_2)$. Then the merged state is
\begin{align}
  \ell &= e^{m_1-m}\ell_1 + e^{m_2-m}\ell_2, \\
  A &= e^{m_1-m}A_1 + e^{m_2-m}A_2.
\end{align}
This follows from the same change-of-reference argument used in the online recurrence. It shows that tiles may be processed independently and merged later so long as their state uses the same sufficient statistics.

\section{Softmax Jacobian contraction}

For one row $p = \softmax(s)$, the Jacobian is
\begin{equation}
  \frac{\partial p_j}{\partial s_t} = p_j(\delta_{jt} - p_t),
\end{equation}
where $\delta_{jt}$ is the Kronecker delta. Let $g_j = \partial \mathcal L/\partial p_j$. Then
\begin{align}
  \frac{\partial \mathcal L}{\partial s_t}
  &= \sum_j g_j \frac{\partial p_j}{\partial s_t} \\
  &= \sum_j g_j p_j (\delta_{jt} - p_t) \\
  &= g_t p_t - p_t \sum_j g_j p_j \\
  &= p_t \left(g_t - \sum_j g_j p_j\right).
\end{align}
In attention, $g_t = G_i v_t^{\mathsf T}$ and the weighted sum becomes $G_i O_i^{\mathsf T}$, which is exactly the scalar $\delta_i$ used by the implementation.

\section{Causal work count}

For square causal attention with sequence length $N$, the number of valid $(i,j)$ pairs is
\begin{equation}
  \sum_{i=1}^{N} i = \frac{N(N+1)}{2}.
\end{equation}
Relative to the dense non-causal case of $N^2$ pairs, the leading ratio tends to one half as $N$ grows:
\begin{equation}
  \frac{N(N+1)/2}{N^2} = \frac{1}{2} + \frac{1}{2N}.
\end{equation}
Finite-size tile boundaries and control flow prevent real runtimes from following this ratio exactly, but it remains a useful leading-order intuition.

\section{Shared-memory footprint}

The forward and gradient kernels stage two FP32 tiles of height $B_c=16$ and width $D$:
\begin{equation}
  2 \times 16 \times D \times 4 \text{ bytes} = 128D \text{ bytes}.
\end{equation}
At $D=64$, $128$, and $256$, this gives 8 KiB, 16 KiB, and 32 KiB respectively. This linear dependence on $D$ is one reason the current native contract stops at $D=256$.

\section{Dense-equivalent FLOP model}

The report uses the common simplified count that treats one multiply-add as two FLOPs.

\paragraph{Forward.}
The score product $QK^{\mathsf T}$ costs approximately $2N_qN_kD$ FLOPs and the value product $PV$ costs another $2N_qN_kD$, giving
\begin{equation}
  F_{\text{fwd}} \approx 4BHN_qN_kD.
\end{equation}

\paragraph{Backward.}
The gradient products for $dQ$, $dK$, and $dV$ contribute the same leading order again, and recomputation adds another score-like pass. The benchmark therefore uses
\begin{equation}
  F_{\text{fwd+bwd}} \approx 12BHN_qN_kD
\end{equation}
as a dense-equivalent estimate. This is a deliberately simple model rather than a complete instruction count.
